\section{Installing the terminal Wilson normalization}
\label{sec:terminal-normalization-transfer}
The preceding startup construction uses a particularly simple terminal
path. We must install it in the common horizon while preserving the
stationary estimates used at the next step. Agreement through fourth
order suffices for the required classical and one-loop coefficients;
it does not identify the full nonlinear stationary maps.

Throughout this section the following data are fixed; they are the
same in every subsection. A paid scale schedule, with stage profiles
whose direct clip coefficients satisfy $0<c_0\le1$. The retained
horizon $M^1$ on that schedule, satisfying the kernel bounds at fixed
kernel prices: a strip width, an upper price and a lower price for
each stage. A local-ratio chart $O$, used in every thermal function
below. The stationary plan
$(R^{\rm old}_{N,m},\alpha^{\rm old}_{N,m},u^{\rm old}_{N,m})$ on
$M^1$ given by Proposition~\ref{prop:path-retained-plan}, planned at
fixed post costs in the sense of Definition~\ref{def:path-plan}. From
that definition we use, at every cutoff $N$ and period: the old
terminal path $R^{\rm old}_{N,N}$ has the Wilson derivatives of every
order $j\le4$ at the identity field and all its loop coefficients
vanish (item 2); and, at the step from $N$ to $N-1$, the gap, the
critical and guessed properties, the classical match and the
finite-cutoff bound (items 3 and 5). Finally the startup path of
Section~\ref{sec:wilson-startup-construction} is assumed to satisfy,
at every stage, the path bounds of that stage at its width and path
allowance, with zero loop allowance; this is arranged by choosing the
stage path allowance above the finite weighted sums of that path, as
explained there.

\subsection{Transporting the terminal quadratic columns}
Write the old terminal full column, at any period, as
\[
 A=\begin{pmatrix}\sqrt{c_0}I\\ R\end{pmatrix},
 \qquad E=\begin{pmatrix}I\\0\end{pmatrix},\qquad Z=A+E.
\]
The terminal metric identity says $A^*A\,d=d$, where $d$ is the
lattice exterior derivative on one-forms. Since $E^*A=A^*E=
\sqrt{c_0}I$, it implies
\[
                  Z^*Ad=Z^*Ed.
\]
Moreover
\[
 Z^*Z=(1+\sqrt{c_0})^2I+R^*R
\]
is positive definite. Put
\[
 P=Z(Z^*Z)^{-1}Z^*,\qquad U=2P-I.
\]
Direct multiplication gives $P^*=P$, $P^2=P$ and $U^*U=I$.
Also $P(Ad+Ed)=Ad+Ed$ and $PAd=PEd$, whence
\[
                            UAd=Ed.
\]
Construct the analogous reflection $U_0$ for the desired terminal
column
$A_0=(\sqrt{c_0}I,\sqrt{1-c_0}I)^{\mathsf T}$.
Then $X=U_0^*U$ is an isometry and $XAd=A_0d$.

These are periodizations of common exponentially decaying kernels.
To check the uniform estimate, suppose the old row kernel has an
exponentially weighted absolute coefficient bound at width $\eta>0$.
Its Fourier symbol and reflected transpose are holomorphic on a
smaller strip. On the real torus the Gram matrix is bounded below
by $(1+\sqrt{c_0})^2I$. The coefficient bound controls uniformly
the change of its symbol under a small imaginary displacement.
Choose a smaller positive width so that this change has operator
norm at most half that lower bound. The Neumann series then
constructs a holomorphic inverse on the whole smaller strip,
bounded by $2/(1+\sqrt{c_0})^2$. Fourier coefficient estimates on
a further smaller strip give an exponentially weighted absolute
sum: shift the contour in each coordinate in the sign of the
Fourier index and sum the resulting geometric series. Matrix
dimension is fixed by the row types, so these constants are
independent of period and block length. Multiplication and
transposition give the same kind of bound for $U,U_0,X$.
This is the strip inverse argument used for the free kernels.

Modify only the finest transfer column $m_f$, replacing it by
$Xm_f$. Since $X^*X=I$, its Gram matrix is unchanged. Its mixed
identity with the terminal curvature column is unchanged after
replacing that column by $A_0$, because $XAd=A_0d$.
The other transfer columns, including the mean column and the
transport matrix, are unchanged. Thus all quadratic transport and
orthogonality identities remain valid. The exponential bounds for
the modified column follow by convolution of its old bound with
the bound for $X$. There is only this one modification at the
terminal step, so its cost does not accumulate along the history.

At each stage choose the new strip width below both the original
width and the width just obtained; choose the new upper price above
one, the old price and this convolution price; and choose the new
lower price below one and the old lower price. These choices depend
only on the stage profiles and old kernel prices. They precede the
block lengths, delays and realization. The resulting horizon agrees
with the old one below its terminal level, its towers agree when
$m+1<N$, and their mean and transport columns agree also at
$m+1=N$. Its terminal direct metric is the identity. This proves
the revised-horizon assertions used in
Section~\ref{sec:common-funded-history}.

\subsection{Finite jets of stationary values}
We record the analytic calculation controlling the change of path.
Let $F,G:X\times Y\to\mathbb R$ be analytic germs at zero on
finite-dimensional real spaces. Suppose their derivatives through
order $k+1$ agree. Let $u,v:X\to Y$ be analytic germs with
$u(0)=v(0)=0$ and
\[
 D_YF(x,u(x))=0,\qquad D_YG(x,v(x))=0.
\]
Assume the vertical Hessian of $F$ at zero is invertible.
Then the graphs $(x,u(x))$ and $(x,v(x))$ agree through order $k$,
and the stationary values $F(x,u(x))$, $G(x,v(x))$ agree through
order $k+1$.

To prove the graph assertion, use the analytic local diffeomorphism
\[
                  \Psi_F(x,y)=(x,D_YF(x,y)).
\]
Its derivative is block triangular with invertible diagonal blocks
$I$ and $D_Y^2F(0)$. The corresponding map $\Psi_G$ has the same
jets through order $k$. On the two stationary graphs these maps
both equal $(x,0)$. Composing their jet identities with the local
inverse of $\Psi_F$ proves equality of graph jets through order
$k$. On a stationary graph the chain rule reduces to
\[
 D\bigl[F(x,u(x))\bigr]=D_XF(x,u(x)),
\]
because the vertical derivative vanishes. The right sides for
$F,u$ and $G,v$ agree through order $k$, using the graph identity
and the derivative jets of $F,G$. Their values at zero also agree.
This proves equality of stationary-value jets through order $k+1$.

In the application $k=3$. The old terminal action and the new
startup action both match the Wilson action through order four.
For the old one this is item (2) of Definition~\ref{def:path-plan},
a property of the old plan which is additional to its quadratic
metric identity; for the new one it is proved in
Section~\ref{sec:wilson-startup-construction}. This matching holds
after substitution of the real fine-field chart on each forest
slice. Indeed that chart is analytic and sends zero to the
identity field; compose the matched real jets with its local
logarithm. The noise term is identical in the two tested actions,
so adding it preserves the matching.

The vertical Hessians at zero are equal, hence the old positive
gap is preserved exactly. The analytic implicit function theorem
constructs the new critical graph. Near zero the tested-action
germs are independent of the branch tag under the chosen kernel's
branch hypothesis; thus this one graph solves every tested
critical equation. Its first derivative agrees with the old graph,
so it has the same linear mean. The preceding calculation also
preserves the classical stationary-value jets through order four.
This is the precise classical matching required by the outgoing
coarse path, which is left unchanged.

\subsection{The one-loop coefficient and retained limits}
The vertical Hessian matrices of the two tested actions have
matching jets through order two, since their actions match through
order four. The determinant is a finite polynomial in matrix
entries, so the determinants have matching two-jets. Composing
with the stationary graphs preserves this matching: the graphs
agree through order three, more than is needed here. Their
positive determinant at zero allows the real logarithm.

The Haar correction is a finite product of the same local chart
Jacobian $J$ along the vertical graph coordinates. Since $J$ is
analytic near zero and $J(0)=1$, these products and their
logarithms also have matching two-jets. Both terminal loop
potentials are zero: the old terminal path has vanishing loop
coefficients by item (2) of Definition~\ref{def:path-plan}, and the
startup path has no loop term. It follows that the two thermal functions
\[
 \frac12\log\det D_Y^2F(x,u(x))
       -\log\!\prod_b J(u(x)_b)
\]
and its counterpart for $G,v$ have identical Hessians at zero.
The finite-period wave estimate is the finite-cutoff bound of
Section~\ref{subsec:path-stationary-chart}: for every forest, every
local choice, every two coarse sites and every two polarization
amplitudes of modulus at most one, the second derivatives at the
origin of the thermal function and of the outgoing comparison
function $\alpha\mathcal K_2+\mathcal L$ of the unchanged coarse
path, on the corresponding plane-wave directions, differ by at most
the stated multiple of $|\operatorname{Site}(c)|e^{-\gamma c}$,
$c$ being the coarse side. It depends only on these Hessians and
the unchanged outgoing classical data. Its bound is therefore
preserved with exactly the old scalar correction, gap and
tolerance. The other two clauses of the finite-cutoff bound are
analyticity of the new thermal function at the origin, which holds
because the new graph, the determinant and the chart Jacobian are
analytic there, and positivity of the determinant, which is the next
remark. Positivity near the new stationary graph follows from
continuity and its positive definite Hessian at zero.

Apply this construction only at $m+1=N$; at every lower step keep
the old graph, path and scalar correction. The new terminal path
is the startup path, while
\[
 R^{\rm new}_{N,m}=R^{\rm old}_{N,m}\quad(m<N),\qquad
 \alpha^{\rm new}_{N,m}=\alpha^{\rm old}_{N,m},
 \qquad u^{\rm new}_{N,m}=u^{\rm old}_{N,m}\quad(m+1<N).
\]
The matching argument above supplies all stationary conditions at
the changed step. The revised quadratic horizon supplies its same
linear guess and all-period kernel identities. For every fixed
level, the changed data consequently agree with the old data on
a tail of cutoffs. Their fixed-level limits and eventual
extended-domain bounds persist by this exact tail equality.
The new terminal graph need not have the old graph's complete
analytic expansion, and no such assertion is used.

\begin{proposition}[Terminal replacement]
\label{prop:terminal-replacement}
With the data fixed at the start of this section, there are revised
kernel prices, depending only on the stage profiles and the old
kernel prices and chosen before the block lengths, a horizon $M$ on
the same schedule satisfying the kernel bounds at the revised
prices, and a stationary plan
$(R_{N,m},\alpha_{N,m},u_{N,m})$ on $M$, planned at the same post
costs and with the same chart $O$ in the sense of
Definition~\ref{def:path-plan}, such that
\begin{align*}
 R_{N,N}&=R^{\rm Wilson}_N,&
 R_{N,m}&=R^{\rm old}_{N,m}\quad(m<N),\\
 \alpha_{N,m}&=\alpha^{\rm old}_{N,m}\quad(m<N),&
 u_{N,m}&=u^{\rm old}_{N,m}\quad(m+1<N),
\end{align*}
where $R^{\rm Wilson}_N$ is the startup path in the clip of level
$N$. The beds satisfy $M.H(N,m)=M^1.H(N,m)$ for $m<N$, and
$M.H(N,N)$ is the bed of the column $A_0$. The towers satisfy
$M.T(N,m)=M^1.T(N,m)$ for $m+1<N$; for $m<N$ their installed columns
and transport columns agree.
\end{proposition}
\begin{proof}
The horizon is constructed in the first subsection, which also gives
the statements on beds and towers and the revised prices. Define the
plan by the displayed identities, with $u_{N,N-1}$ the new critical
graph of the second subsection. We check
Definition~\ref{def:path-plan}. Item (1): for $m<N$ the paths and
beds are the old ones; the startup path is attached to the terminal
bed of $A_0$ because its only row entry is $\sqrt{1-c_0}$ at offset
zero, and it satisfies the stage bounds by the standing assumption.
Item (2) is Section~\ref{sec:wilson-startup-construction}. Items (3)
and (5) are unchanged for $m+1<N$, since the incoming and outgoing
paths, the tower and the maps are the old ones. For $m+1=N$ they are
the gap, critical, guessed and classical statements of the second
subsection and the thermal statement of the third, all for the
unchanged outgoing path $R^{\rm old}_{N,N-1}$ and scalar
$\alpha^{\rm old}_{N,N-1}$; the installed column used in the guessed
property is unchanged. Item (4) involves only $\alpha$ and the
schedule.
\end{proof}

This is the replacement used in
Section~\ref{sec:common-funded-history}. It selects no new background
and no new cutoff sequence.
